\section{Experimental Setup}
\label{sec:experimental_setup}

\subsection{Datasets}

We evaluate on four sources of ambiguity-injected coding problems:

\begin{itemize}
\item \textbf{ClarifyCodeBench}~\cite{clarifycodebench}: 419 problems
derived from LiveCodeBench, each annotated with $K \in \{1,2,3\}$
deliberately deleted specification details.
\item \textbf{Orchid-HEval}: HumanEval~\cite{humaneval} problems rewritten
with a single injected ambiguous edit (lexical, vagueness, semantic, or
syntactic), paired with the original hidden tests.
\item \textbf{Orchid-BCB} and \textbf{Orchid-BCB-Expand}: BigCodeBench
\cite{bigcodebench} problems rewritten the same way, with BCB-Expand a
larger, semi-automatically constructed extension of the original BCB
subset.
\end{itemize}

All source-problem groupings are preserved under splitting: no two prompt
variants derived from the same underlying problem appear on opposite sides
of a train/validation/test or leave-one-source-out (LODO) split. We use
grouped, stratified splits throughout the ambiguity-degree probe
evaluation, and full leave-one-source-out evaluation (train on three
sources, test on the fourth, held out entirely) for the ranking probe, to
measure whether ranking quality generalizes across benchmark construction
methodologies rather than only within one.

\subsection{Models}

The probe-host model is Qwen3-8B (bf16, thinking mode disabled) throughout.
For code generation, we use two models of different capability: Qwen3-8B
itself, and DeepSeek V3.2 (accessed via Bedrock), to test whether findings
about clarification quality are specific to a particular coding model or
hold across models of different capability. Candidate clarifying-question
elicitation and independent quality judging (assigning tiers 0/1/2 to
candidates relative to the benchmark's labeled ambiguity) use DeepSeek V3.2
as an independent judge, separate from whichever model generates the code
under test.

\subsection{Baselines}

\textbf{Zero-shot self-rank}: the coding model is prompted with the
ambiguous prompt and the full numbered list of candidate clarifying
questions in a single message, and asked to output the number of the one it
would ask. This requires exactly one generation call per problem (greedy
decoding, no sampling), with the entire candidate pool inlined -- not one
call per candidate. \textbf{No-clarification}: the raw ambiguous prompt is
passed directly to the coding model with no clarification step at all,
establishing the baseline value of clarifying at all, independent of which
question is chosen.

\subsection{Correctness Metric}

We report correctness as the \textbf{fractional test-case pass rate}: the
number of individual test assertions a generated solution passes divided by
the total number of test assertions for that problem, averaged across
groups. This is distinct from -- and, as shown in
Section~\ref{sec:results}, statistically more sensitive than -- an
all-or-nothing metric that scores a group as passing only if every test
assertion succeeds. Both metrics are computable from the same underlying
per-group pass/total counts; we report the fractional metric as primary
throughout and note where it diverges from the all-or-nothing metric's
conclusions.

\subsection{Significance Testing}

For between-groups comparisons (e.g., groups where the probe ranked the
labeled ambiguity first vs.\ groups where it did not), we report
Mann-Whitney U and Welch's $t$-test alongside a 10,000-resample bootstrap
95\% confidence interval on the mean difference. For within-groups,
paired comparisons (e.g., the same 244 problems with vs.\ without any
clarification), we report the Wilcoxon signed-rank test, a paired $t$-test,
and a paired bootstrap 95\% CI on the mean of per-group differences. We
disclose cases where these tests disagree at the margin rather than
reporting only the test that reaches $p < 0.05$.

\subsection{Compute Environment}

\label{sec:experimental_setup:compute}
All hidden-state extraction and generation runs on a single GPU
(NVIDIA A10G, 23\,GB, native bf16, no CPU weight offload). We note that an
earlier measurement on a memory-constrained T4 GPU requiring CPU offload
showed no measurable speedup between the probe and zero-shot generation, an
artifact of a fixed CPU$\leftrightarrow$GPU weight-transfer cost dominating
both paths equally; we report only native-GPU timings here, and describe
the exact benchmarked comparison in
Section~\ref{sec:results:cost}.
